% %%%%%%%%%%%%%%%%%%%%%%%%%%%%%%%%%%%%
%   Dependencies
% %%%%%%%%%%%%%%%%%%%%%%%%%%%%%%%%%%%%
\usepackage[T1]{fontenc}
\usepackage[export]{adjustbox}
\usepackage{tabularx}
\usepackage{longtable}
\usepackage{graphicx}
\usepackage{colortbl}
\usepackage{lastpage}
\usepackage[T1]{fontenc}
\usepackage[utf8]{inputenc}
\usepackage{eurosym}
\usepackage{multirow}
\usepackage{ulem}
\usepackage{soul}
\usepackage{fancyhdr}
\usepackage{geometry}
\usepackage{indentfirst}
\usepackage{subcaption}
\usepackage{float}
\usepackage[default,oldstyle,scale=0.95]{opensans}
\usepackage{xcolor}
\usepackage{comment}
\usepackage{amsmath}
\usepackage{enumitem}
\usepackage[font=small,labelfont=bf]{caption}
\usepackage{subfigure} % subfiguras
\usepackage{float} % Podemos usar o especificador [H] nas figuras para que fiquem na posição desejada
\usepackage{capt-of} % Permite usar etiquetas fora de elementos flutuantes
% (etiquetas de figuras)
\usepackage[colorlinks=true,plainpages=true,citecolor=blue,linkcolor=blue]{hyperref} % Para que os hypertextos apareçam coloridos
\usepackage{listings} % Para usar código fonte
	\definecolor{dkgreen}{rgb}{0,0.6,0} % Definimos cores para usar no código
	\definecolor{gray}{rgb}{0.5,0.5,0.5} 
% configuração para a linguagem de programação que desejamos utilizar
	\lstset{language=Matlab,
   	keywords={break,case,catch,continue,else,elseif,end,for,function,
    	  global,if,otherwise,persistent,return,switch,try,while},
   	basicstyle=\ttfamily,
   	keywordstyle=\color{blue},
   	commentstyle=\color{red},
   	stringstyle=\color{dkgreen},
   	numbers=left,
   	numberstyle=\tiny\color{gray},
   	stepnumber=1,
   	numbersep=10pt,
   	backgroundcolor=\color{white},
   	tabsize=4,
   	showspaces=false,
   	showstringspaces=false}
\DeclareGraphicsExtensions{.pdf,.png,.jpg}
\graphicspath{ {graphics/} }
